\documentclass[12pt]{article}
\usepackage[utf8]{inputenc}
\usepackage[T1]{fontenc}
\usepackage{hyperref}
\title{Kap Language Resources for K-12 Curriculum Development}
\author{Compiled by Claude Sonnet 5.5 (claude-sonnet-5-5) for Andr\'as Kornai}
\date{}
\begin{document}
\maketitle

\section{Introduction}
Kap (ISO 639-3 ykm, Glottocode kapp1237) is an Austronesian language of Sandaun Province, Papua New Guinea. Search results also list it as Yakamul, with Ali, Yakamul and Yigel given as varieties or names. It is spoken in East Aitape Rural LLG, including the village of Yakamul and the islands of Ali, Angel and Seleo. Speaker estimates vary widely: about 3,500 (Wikipedia, 2003), 2,079 to 3,460 (Verbix) and 6,500 (Joshua Project). Joshua Project reports a shift to Tok Pisin and says most children do not speak Kap. The search found no school materials in Kap.

\section{Primary Educational Resources}
The searches found no primers, textbooks or school programmes in Kap.

\section{Linguistic and Reference Materials}
\begin{itemize}
\item Wikipedia, ``Yakamul language'': overview of classification, location and speaker numbers. \url{https://en.wikipedia.org/wiki/Yakamul_language}
\item Verbix, Kap language page (an endangered-language page also exists at \url{https://docs.verbix.com/EndangeredLanguages/Kap}). \url{https://docs.verbix.com/Languages/Kap}
\item Endangered Languages Project entry. \url{https://endangeredlanguages.com/lang/10522?hl=ru}
\end{itemize}
I saw no grammar or dictionary in the search results.

\section{Religious and Community Resources}
\begin{itemize}
\item Joshua Project, Yakamul in Papua New Guinea: people-group profile covering language shift and the lack of Kap-speaking children. \url{https://joshuaproject.net/people_groups/15955/PP}
\end{itemize}
I found no community organisations.

\section{Digital and Audio Resources}
\begin{itemize}
\item Global Recordings Network, language page for ykm. The search summary said no recordings are currently available at the ISO level; recordings in some Kap varieties may be listed under another entry. \url{https://www.globalrecordings.net/en/language/ykm}
\end{itemize}

\section{Curriculum Development Context}
Very little documented material exists for Kap. A curriculum would have to start from new documentation: wordlists, a practical orthography and community-recorded stories. Children are shifting to Tok Pisin, so early-years use of Kap and community involvement matter. English and Tok Pisin materials would supply the rest of the curriculum for now.

\end{document}